\documentclass[12pt]{article}
% ---------- Packages ----------
\usepackage[utf8]{inputenc}
\usepackage[T1]{fontenc}
\usepackage{lmodern}
\usepackage[margin=1in]{geometry}
\usepackage{microtype}
\usepackage{mathtools, amsmath, amssymb, amsthm}
\usepackage{bbm}
\usepackage{enumitem}
\usepackage{hyperref}
\usepackage[nameinlink,capitalize]{cleveref}
\usepackage{csquotes}

\title{On Vanishing Sums and Uniform Bounds in the Theta--Mellin Determinant Framework}
  

\begin{document}
\maketitle
\begin{center}

\LARGE{Citizen Gardens} 
\vspace{0.5cm} \\
\textit{Institute for Mathematical Discovery} \\
\vspace{0.5cm}
{\normalsize \today \par}
\end{center}
\vspace{3cm}  
\section{Vanishing of Sums at Special Imaginary Values}

Within the Theta--Mellin framework, the central identity proved in the cited work is
\begin{equation}
\det\nolimits_{2}\big(I - X(s)\big) = C(s)\, \xi(s),
\end{equation}
where $C(s)$ is an entire, nowhere-vanishing function satisfying $C(1-s) = C(s)$. Thus the determinant vanishes if and only if $\xi(s)=0$.

\subsection{Analytic Mechanism}
The collapse of sums or integrals to zero occurs at those imaginary parts corresponding to nontrivial zeros of $\xi(s)$. Since $C(s)\neq 0$ everywhere, the only way for the determinant (or any equivalent regularized sum) to vanish is when $\xi(s)=0$.

\subsection{Theta--Mellin Cancellation}
The Mellin identity for the Jacobi theta tail,
\begin{equation}
\int_{0}^{\infty} (\Theta(u)-1) \, u^{z/2 - 1} \, du = \frac{2\, \xi(z)}{z(z-1)},
\end{equation}
implies that if $z$ is a zero of $\xi$, then the integral vanishes. At such frequencies, the oscillatory weight $u^{i\,\mathrm{Im}(z)/2}$ forces exact cancellation.

\subsection{Spectral Interpretation}
Under the Mellin transform, $X(s)$ diagonalizes with symbol
\begin{equation}
m_s(t) = \frac{\xi\left(\tfrac{s}{2}+it\right)}{\left(\tfrac{s}{2}+it\right)\left(\tfrac{s}{2}+it-1\right)}.
\end{equation}
The regularized log-determinant is
\begin{equation}
\log \det\nolimits_{2}(I - X(s)) = \int_{\mathbb{R}} \frac{\log(1-m_s(t)) - m_s(t)}{2\pi}\, dt.
\end{equation}
Near a simple zero $s=\rho$ of $\xi$, the integrand reproduces the singularity of $\log \xi(s)$, forcing $\det_2(I-X(s))=0$ precisely at $s=\rho$.

\section{Uniform $\varepsilon$--Bounds Away from Zeros}

Since $C(s)$ never vanishes, one can show that away from the set of zeros of $\xi$, the determinant remains bounded away from zero.

\subsection{Global Compact Bound}
Let $K\subset \mathbb{C}$ be compact and define the $\delta$--exclusion of the zero set
\begin{equation}
K_{\delta} := \{ s\in K : \operatorname{dist}(s, \mathcal{Z}(\xi)) \geq \delta \}.
\end{equation}
Then the continuous function $s \mapsto |C(s)\,\xi(s)|$ attains a positive minimum on $K_{\delta}$:
\begin{equation}
\varepsilon(K,\delta) := \inf_{s\in K_{\delta}} |\det_{2}(I - X(s))| > 0.
\end{equation}
Thus any deviation of at least $\delta$ from the zero set yields a uniform positive lower bound.

\subsection{Local Bound Near a Zero}
Suppose $\rho$ is a zero of multiplicity $m\geq 1$. Then locally,
\begin{equation}
\xi(s) = a_m (s-\rho)^m (1+o(1)), \qquad a_m\neq 0.
\end{equation}
For sufficiently small $r$, there exists $c_\rho>0$ such that
\begin{equation}
\min_{|s-\rho|=r} |\xi(s)| \geq c_\rho \, r^m.
\end{equation}
Since $C(s)$ is continuous and nonvanishing, $\min_{|s-\rho|=r}|C(s)|\geq c_C>0$. Hence,
\begin{equation}
\min_{|s-\rho|=r} |\det_{2}(I - X(s))| \geq (c_C c_\rho) r^m =: \varepsilon_\rho(r) > 0.
\end{equation}
Thus exact vanishing occurs only at $s=\rho$; for any fixed nonzero deviation $r$ from $\rho$, the determinant is bounded below by a positive power law.

\section{Conclusion}
Specific imaginary parts collapse the theta--Mellin sum precisely when they correspond to ordinates of nontrivial zeros of $\xi$. Away from those points, one obtains explicit $\varepsilon$--bounds: globally uniform on compact sets excluding neighborhoods of zeros, and locally quantified in terms of the zero multiplicity.

\end{document}
